\chapter{向量空间}\label{ch:b1:vspaces}

线性代数从这里开始：\omterm{def:b1:vspaces:def}{向量空间}的公理抽取出
$\R^2$、$\R^3$、\omterm{def:b1:poly:def}{多项式}空间与函数空间的共同点
——即可以相加，也可以作数乘。两章建立了这一理论
（\cref{ch:b1:findim} 引入维数）；它们建立的语言——
\omterm{def:b1:vspaces:span}{张成的子空间}、\omterm{def:b1:vspaces:free}{线性无关}族、基、直和——是此后每一章都离不开的
工具。全文中，$K$ 表示 $\R$ 或 $\C$
（\emph{标量}）。

\section{定义与例子}

\begin{definition}[向量空间]\label{def:b1:vspaces:def}
称 \emph{$K$-向量空间}\index{向量空间}为\omterm{def:b1:logic:sets}{集合} $E$，
在其上定义加法，使 $(E, +)$ 构成\omterm{def:b1:structures:group}{阿贝尔群}（零元记作 $0_E$ 或
$0$），并定义数乘 $K \times E \to E$，使对任意
$\lambda, \mu \in K$ 与 $x, y \in E$，下列等式成立：
\[
\lambda(x + y) = \lambda x + \lambda y,\quad
(\lambda + \mu) x = \lambda x + \mu x,\quad
\lambda(\mu x) = (\lambda\mu) x,\quad
1\,x = x .
\]
推论：$0\,x = 0_E$，$\lambda\,0_E = 0_E$，$(-1)x = -x$，以及
$\lambda x = 0_E \implies \lambda = 0$ 或 $x = 0_E$（两边乘以
$\lambda^{-1}$）。
\end{definition}

\begin{proof}[推论的证明]
对 $0\,x = 0_E$：由 $(0 + 0)x = 0x + 0x$ 与 $(0+0)x = 0x$，
消去 $0x$（运算在群 $(E, +)$ 中进行）。对 $\lambda\,0_E$，同样的技巧用于
$\lambda(0_E + 0_E)$。对 $(-1)x$，加上 $x$，
\[
x + (-1)x = 1\,x + (-1)x = \bigl(1 + (-1)\bigr)x = 0\,x = 0_E ,
\]
于是 $(-1)x$ 即为 $x$ 的加法逆元。最后，若 $\lambda x
= 0_E$ 且 $\lambda \neq 0$，两边乘以 $\lambda^{-1}$（
标量构成一个域），并利用以下两条公理
$\lambda^{-1}(\lambda x) = (\lambda^{-1}\lambda)x = 1x = x$
与 $\lambda^{-1}0_E = 0_E$ 一起，得 $x = 0_E$。这四条规则虽简单，
却在此后的每一页上被不加说明地用到——
其中最后一条恰恰说明为何需要域：
若标量取 $\Z$，“空间” $\Z/2\Z$ 便会
在 $2\,x = 0$ 处违反它。
\end{proof}

\begin{example}\label{ex:b1:vspaces:examples}
$K^n$（按分量运算）；\omterm{def:b1:poly:def}{多项式}全体 $K[X]$；
$\mathcal{F}(A, K)$ 即从任意\omterm{def:b1:logic:sets}{集合} $A$ 到 $K$ 的函数全体（按点运算）——
包括\omterm{def:b1:continuity:continuous}{连续}函数、数列
$\mathcal{F}(\N, \R)$ 等；$\C$ 作为 $\R$ 上的\omterm{def:b1:vspaces:def}{向量空间}。在每种
情形下，公理均自 $K$ 继承而来。
\end{example}

\begin{definition}[子空间]\label{def:b1:vspaces:subspace}
称 $F \subseteq E$ 为\emph{子空间}\index{子空间}，若 $0_E \in
F$，且 $F$ 对加法及数乘封闭——
等价地：
\[
F \neq \emptyset
\qquad\text{且}\qquad
\forall x, y \in F,\ \forall \lambda \in K,\quad
x + \lambda y \in F .
\]
子空间本身也是\omterm{def:b1:vspaces:def}{向量空间}。任意多个子空间的交仍是子空间；
而并集几乎从不是子空间（证明与
\cref{exo:b1:structures:6} 相同）。
\end{definition}

\begin{example}\label{ex:b1:vspaces:subspaces}
在 $\mathcal{F}(\R, \R)$ 中：\omterm{def:b1:continuity:continuous}{连续}函数、
\omterm{def:b1:derivative:def}{可导}函数、次数 $\leq n$ 的\omterm{def:b1:poly:def}{多项式}（记作
$K_n[X]$，在 $K[X]$ 中）、齐次线性
微分方程的解（\cref{thm:b1:diffeq:homogeneous2} 已言明此事）。
反例：$\{f : f(0) = 1\}$（不含零元），次数恰为
$n$ 的\omterm{def:b1:poly:def}{多项式}（对加法不封闭）。
\end{example}

\begin{example}[是否为子空间：四个结论及论证]
\label{ex:b1:vspaces:subspaceornot}
在实数数列空间中：
\begin{itemize}
  \item $\{u : u \text{有界}\}$ \emph{是}\omterm{def:b1:vspaces:subspace}{子空间}：$0$
        有界，若 $\abs{u_n} \leq M$，$\abs{v_n} \leq M'$，
        则 $\abs{u_n + \lambda v_n} \leq M + \abs\lambda M'$。
  \item $\{u : u_n \to 1\}$ \emph{不}是\omterm{def:b1:vspaces:subspace}{子空间}：零数列
        不在其中（且两元素之和趋于 $2$）。
  \item $\{u : u \text{单调}\}$ \emph{不}是\omterm{def:b1:vspaces:subspace}{子空间}：$u_n = n$
        与 $v_n = -n + (-1)^n$ 都单调，它们的和
        $(-1)^n$ 不单调；失效的正是对加法的封闭性，
        尽管该\omterm{def:b1:logic:sets}{集合}含有 $0$，以及
        其元素的全部数乘。
  \item $\{u : u_{n+1} = u_n^2\}$ \emph{不}是\omterm{def:b1:vspaces:subspace}{子空间}：它含有 $0$，
        但 $2u$ 在 $u$ 为该\omterm{def:b1:logic:sets}{集合}中的非零元时便脱离该\omterm{def:b1:logic:sets}{集合}
        （一般地 $2u_{n+1} \neq (2u_n)^2$）——平方运算
        正是非线性的所在。
\end{itemize}
判断时的步骤总是相同：先检验 $0$（最省事），
再检验封闭性——要否定某个结论时，
一对明确的反例胜过任何疑虑。
\end{example}

\section{张成、和与直和}

\begin{definition}[线性组合、张成的子空间]\label{def:b1:vspaces:span}
\emph{线性组合}是指向量族 $(x_1, \dots, x_p)$（$E$ 中的向量）的
任一形如 $\lambda_1 x_1 + \dots + \lambda_p x_p$ 的向量
（$\lambda_i \in K$）。全体线性组合构成的\omterm{def:b1:logic:sets}{集合}称为
\emph{张成的子空间}\index{张成的子空间} $\operatorname{Vect}(x_1, \dots, x_p)$：它
是\omterm{def:b1:vspaces:subspace}{子空间}，且是包含该族的最小\omterm{def:b1:vspaces:subspace}{子空间}。
\end{definition}

\begin{proof}[两个断言的证明]
稳定性：两个线性组合之和 $\sum\lambda_i x_i +
\sum\mu_i x_i = \sum(\lambda_i + \mu_i)x_i$ 仍为线性组合，
数乘亦然：$\mu\sum\lambda_i x_i = \sum(\mu\lambda_i)
x_i$；零组合表明 $0$ 属于其中，于是\omterm{def:b1:vspaces:span}{张成的
子空间}。极小性：设 $H$ 为任一包含 $x_1,
\dots, x_p$ 的\omterm{def:b1:vspaces:subspace}{子空间}。由数乘下的稳定性，每个
$\lambda_i x_i \in H$，由加法下的稳定性，它们的和
落在 $H$ 中：每个线性组合都属于 $H$，即\
$\operatorname{Vect}(x_1, \dots, x_p) \subseteq H$。因此\omterm{def:b1:vspaces:span}{张成的
子空间}包含于每一个包含该族的\omterm{def:b1:vspaces:subspace}{子空间}，即它是
最小的\omterm{def:b1:vspaces:subspace}{子空间}。
\end{proof}

\begin{definition}[和、直和]\label{def:b1:vspaces:sum}
设 $F, G$ 为 $E$ 的\omterm{def:b1:vspaces:subspace}{子空间}：
\[
F + G = \{\,u + v : u \in F,\ v \in G\,\}
\]
是\omterm{def:b1:vspaces:subspace}{子空间}（包含 $F \cup G$ 的最小\omterm{def:b1:vspaces:subspace}{子空间}）。该和
为\emph{直和}\index{直和}，记作 $F \oplus G$，条件是 $F + G$ 中
每个元素都可\emph{唯一}地分解为 $u + v$；
等价地（见下文）即 $F \cap G = \{0\}$。当 $E = F \oplus
G$ 时，称两\omterm{def:b1:vspaces:subspace}{子空间}为\emph{补子空间}\index{补子空间}
在 $E$ 中。
\end{definition}

\begin{example}[两直线之和]
\label{ex:b1:vspaces:sumoflines}
在 $\R^3$ 中，设 $F = \operatorname{Vect}\bigl((1,0,1)\bigr)$ 与
$G = \operatorname{Vect}\bigl((0,1,1)\bigr)$。它们的和为
\[
F + G = \{\,a(1,0,1) + b(0,1,1)\,\}
= \{(a,\ b,\ a + b)\} = \{(x, y, z) : z = x + y\},
\]
即过原点且包含两直线的平面。它
严格大于并集 $F \cup G$（仅是两直线相交而成的十字形）。
例如向量 $(1, 1, 2) = (1,0,1) + (0,1,1)$ 属于和，却不在任一直线上。
又 $F \cap G = \{0\}$（公共向量须满足
$a(1,0,1) = b(0,1,1)$，其前两个\omterm{prop:b1:vspaces:coordinates}{坐标}迫使
$a = b = 0$），故和为直和，且 $F \oplus
G$ 恰为该平面。
\end{example}

\begin{proposition}\label{prop:b1:vspaces:directsum}
$F + G$ 为直和，当且仅当 $F \cap G = \{0\}$。
\end{proposition}

\begin{proof}
若某 $w \neq 0$ 属于 $F \cap G$：$w = w + 0 = 0 + w$ 是 $w$ 的两种
分解。反之，若 $u + v = u' + v'$，其中 $u, u'
\in F$，$v, v' \in G$，则 $u - u' = v' - v$ 属于 $F \cap G =
\{0\}$，故分解唯一。
\end{proof}

\begin{method}[证明 $E = F \oplus G$]
\label{met:b1:vspaces:directsum}
需检验两点，各有其标准的入手方式。
\begin{enumerate}
  \item \emph{交集平凡。}取 $x \in F \cap G$，
        写出两个属于条件，并由此推出 $x =
        0$（切勿凭图形论证，参见下文易错点）。
  \item \emph{和为全空间。}任取 $x \in E$
        并\emph{构造}分解 $x = f + g$——
        或由目标猜出 $f$（$f$ 须满足
        $F$ 的定义性质，它通常决定其公式），
        或解线性方程组，将 $x$
        用 $F$、$G$ 的生成元表示。
\end{enumerate}
若分解式是猜出的，则唯一性
由第 1 步自动得到；若只是存在性不明，则
主要工作在第 2 步。下面两例演示此方法：对于
偶、奇函数，$f$ 的公式由在 $x$ 与 $-x$ 处代入所需恒等式确定；
对于在某点取零的\omterm{def:b1:poly:def}{多项式}，
则由在 $a$ 处求值确定。
\end{method}

\begin{example}\label{ex:b1:vspaces:evenodd}
在 $\mathcal{F}(\R, \R)$ 中，偶函数集 $\mathcal{P}$ 与
奇函数集 $\mathcal{I}$ 互为\omterm{def:b1:vspaces:sum}{补子空间}：任一 $f$ 可写成
\[
f(x) = \underbrace{\frac{f(x) + f(-x)}{2}}_{\text{偶}}
+ \underbrace{\frac{f(x) - f(-x)}{2}}_{\text{奇}},
\]
且既偶又奇的函数恒为零。（取 $\exp$ 时，此分解即
在\cref{ch:b1:functions}中的 $(\cosh, \sinh)$ 这一对。）
\end{example}

\begin{example}[{$K_n[X]$ 中的一对补子空间}]
\label{ex:b1:vspaces:hyperplane}
固定 $a \in K$，令 $F = \{P \in K_n[X] : P(a) = 0\}$，$G =
\operatorname{Vect}(1)$（常值\omterm{def:b1:poly:def}{多项式}）。于是 $K_n[X] = F \oplus
G$。事实上 $F \cap G$ 即在 $a$ 处取零的常值\omterm{def:b1:poly:def}{多项式}，
即\ $\{0\}$；且每个 $P$ 都可分解为
\[
P = \underbrace{\bigl(P - P(a)\bigr)}_{\in F}
+ \underbrace{P(a)}_{\in G} .
\]
此分解值得牢记：减去某点处的值，
是向“在
$a$ 处取零的函数”作投影的标准方法。注意 $F$ 是大\omterm{def:b1:vspaces:subspace}{子空间}，$G$ 是小\omterm{def:b1:vspaces:subspace}{子空间}；一
对\omterm{def:b1:vspaces:sum}{补子空间}未必在任何意义上均衡。
\end{example}

\begin{example}[补子空间从不唯一]
\label{ex:b1:vspaces:manysupplements}
在 $\R^2$ 中，设 $F = \operatorname{Vect}\bigl((1,0)\bigr)$（
$x$ 轴）。$G = \operatorname{Vect}\bigl((0,1)\bigr)$ 与
$G' = \operatorname{Vect}\bigl((1,1)\bigr)$ 均为
$F$ 的\omterm{def:b1:vspaces:sum}{补子空间}：每个都只与 $F$ 相交于 $0$，且每对之和为 $\R^2$。
同一向量的分解各不相同：
\[
(2,\ 1.5) = \underbrace{(2, 0)}_{\in F} +
\underbrace{(0, 1.5)}_{\in G}
= \underbrace{(0.5,\ 0)}_{\in F} +
\underbrace{(1.5,\ 1.5)}_{\in G'} .
\]
事实上，除 $F$ 自身外，\emph{每一条}直线都是
$F$ 在 $\R^2$ 中的\omterm{def:b1:vspaces:sum}{补子空间}：\omterm{def:b1:vspaces:sum}{补子空间}很多，
谈论“那一个”\omterm{def:b1:vspaces:sum}{补子空间}毫无意义，除非
引入额外结构（内积，见\cref{ch:b1:euclid}），
才能选出其中一个。
\end{example}

\begin{omfigure}
\begin{tikzpicture}[scale=1.5]
  \draw[->, thin] (-0.4,0) -- (2.9,0) node[below] {$x$};
  \draw[->, thin] (0,-0.4) -- (0,1.9) node[left] {$y$};
  \draw[omDef, very thick] (-0.3,0) -- (2.7,0)
    node[below, pos=0.95, font=\small] {$F$};
  \draw[omThm, very thick] (0,-0.3) -- (0,1.7)
    node[left, pos=0.95, font=\small] {$G$};
  \draw[omProp, very thick] (-0.25,-0.25) -- (1.75,1.75)
    node[right, pos=0.97, font=\small] {$G'$};
  \fill (2,1.5) circle (0.045) node[above right, font=\small]
    {$(2, 1.5)$};
  \draw[omThm, dashed] (2,1.5) -- (2,0);
  \draw[omDef, dashed] (2,1.5) -- (0,1.5);
  \draw[omProp, dashed] (2,1.5) -- (0.5,0);
  \fill (2,0) circle (0.035);
  \fill (0.5,0) circle (0.035);
\end{tikzpicture}

{\small 同一点在 $\R^2$ 中沿 $F$ 的两种分解
（$x$ 轴）：取\omterm{def:b1:vspaces:sum}{补子空间} $G$（竖直下降）与
\omterm{def:b1:vspaces:sum}{补子空间} $G'$（斜向下降）。两种分解中 $F$ 分量
不同：投影依赖于下降的方向。}
\end{omfigure}

\section{线性无关族、生成族、基}

\begin{definition}\label{def:b1:vspaces:free}
设 $(x_1, \dots, x_p)$ 为 $E$ 中的向量族，定义如下：
\begin{itemize}
  \item 当该族\emph{生成}（$E$）时，即
        $\operatorname{Vect}(x_1,\dots,x_p) = E$；
  \item \emph{线性无关}\index{线性无关}（即其
        向量\emph{线性无关}）当
        \[
        \lambda_1 x_1 + \dots + \lambda_p x_p = 0
        \implies \lambda_1 = \dots = \lambda_p = 0 ;
        \]
        否则称该族\emph{线性相关}；
  \item 当其既线性无关又生成时，称为\emph{基}\index{基}。
\end{itemize}
\end{definition}

\begin{proposition}[坐标]\label{prop:b1:vspaces:coordinates}
$(e_1, \dots, e_n)$ 是 $E$ 的基，当且仅当每个 $x \in E$
都可以\emph{唯一}地写成 $x = \lambda_1 e_1 + \dots +
\lambda_n e_n$ 的形式；标量 $\lambda_i$ 称为 $x$ 在该基下的
\emph{坐标}\index{坐标}。
\end{proposition}

\begin{proof}
生成性 $=$ 分解的存在性。唯一性 $=$
\omterm{def:b1:vspaces:free}{线性无关}性：同一 $x$ 的两种分解之差是
一个等于 $0$ 的组合；由\omterm{def:b1:vspaces:free}{线性无关}性，其全部系数——
即\omterm{prop:b1:vspaces:coordinates}{坐标}之差——必为零。反之，一个非平凡的
零组合给出两种分解 $0 = \sum \lambda_i
e_i = \sum 0\,e_i$.
\end{proof}

\begin{example}\label{ex:b1:vspaces:bases}
$K^n$ 的\emph{标准基}：$e_i = (0, \dots, 1, \dots, 0)$
（$1$ 位于第 $i$ 个位置）。单项式 $(1, X, X^2, \dots, X^n)$ 是
$K_n[X]$ 的基（\omterm{def:b1:vspaces:free}{线性无关}性：零组合即零\omterm{def:b1:poly:def}{多项式}，所以
全部系数为零，见\cref{def:b1:poly:def}）。
在 $\C$ 看作 $\R$ 上的\omterm{def:b1:vspaces:def}{向量空间}时，取基 $(1, \iu)$。
\end{example}

\begin{remark}[坐标是集体协作的结果]
\label{rem:b1:vspaces:teameffort}
$x$ 在基 $(e_1, \dots, e_n)$ 下的第一个\omterm{prop:b1:vspaces:coordinates}{坐标}
取决于\emph{全部}基向量，而不只是 $e_1$。在
$\R^2$ 中，向量 $(3, 1)$ 在标准基下的第一个\omterm{prop:b1:vspaces:coordinates}{坐标}为 $3$，
而在另一组基中，其第一个\omterm{prop:b1:vspaces:coordinates}{坐标}为 $2$，
这组基为 $\bigl((1,0), (1,1)\bigr)$——解方程 $(3,1) = a(1,0) + b(1,1)$：
$b = 1$，$a = 2$。改变一个基向量会重新分配
\emph{所有}\omterm{prop:b1:vspaces:coordinates}{坐标}；\cref{ch:b1:matrices} 将把这种
重新分配归纳为过渡矩阵。
\end{remark}

\begin{example}[检验候选基，从头到尾]
\label{ex:b1:vspaces:testbasis}
$\mathcal{F} = (1 + X,\ 1 + X^2,\ X + X^2)$ 是否为
$\R_2[X]$ 的基？记这三个\omterm{def:b1:poly:def}{多项式}为 $u_1, u_2, u_3$。
\emph{\omterm{def:b1:vspaces:free}{线性无关}性}：零组合 $a\,u_1 + b\,u_2 + c\,u_3 =
0$ 给出，逐系数比较：
\[
a + b = 0, \qquad a + c = 0, \qquad b + c = 0 ;
\]
两式相减得 $b = c$，再由第三式得 $2b =
0$：$a = b = c = 0$，即\omterm{def:b1:vspaces:free}{线性无关}。\emph{生成性}：无需
解三个方程组，注意到如下对称组合
\[
u_1 + u_2 - u_3 = (1 + X) + (1 + X^2) - (X + X^2) = 2 ,
\]
于是 $1 = \frac12(u_1 + u_2 - u_3)$；然后
\[
X = u_1 - 1 = \tfrac12\bigl(u_1 - u_2 + u_3\bigr),
\qquad
X^2 = u_2 - 1 = \tfrac12\bigl(-u_1 + u_2 + u_3\bigr).
\]
单项式都落在由 $\mathcal{F}$ \omterm{def:b1:vspaces:span}{张成的子空间}中，于是一切\omterm{def:b1:poly:def}{多项式}亦然，
故它是基。作为额外收获，合并上面三式，可得
任意 $P = \alpha + \beta X + \gamma X^2$ 的\omterm{prop:b1:vspaces:coordinates}{坐标}：
\[
P = \frac{\alpha + \beta - \gamma}{2}\,u_1
+ \frac{\alpha - \beta + \gamma}{2}\,u_2
+ \frac{-\alpha + \beta + \gamma}{2}\,u_3 .
\]
（验算：取 $P = X$，其\omterm{prop:b1:vspaces:coordinates}{坐标}为 $\bigl(\frac12,
-\frac12, \frac12\bigr)$，与上文所求一致）。两点启示：
族中的对称性往往隐藏着一个捷径组合；而一旦
维数可用（见\cref{ch:b1:findim}），整个
关于生成性的那一半工作将不费吹灰之力——
三个\omterm{def:b1:vspaces:free}{线性无关}向量在 $3$ 维空间中必构成基。
\end{example}

\begin{proposition}[实用的线性无关判别准则]\label{prop:b1:vspaces:freecriteria}
\begin{enumerate}
  \item 由\emph{两两次数互异的
        非零\omterm{def:b1:poly:def}{多项式}}组成的族\omterm{def:b1:vspaces:free}{线性无关}。
  \item 向\omterm{def:b1:vspaces:free}{线性无关}族中添加一个向量，所得族仍\omterm{def:b1:vspaces:free}{线性无关}，当且仅当
        该向量位于此族\omterm{def:b1:vspaces:span}{张成的子空间}之外。
  \item \omterm{def:b1:vspaces:free}{线性无关}族的任一子族仍\omterm{def:b1:vspaces:free}{线性无关}；
        包含生成族的任一族都是生成族。
\end{enumerate}
\end{proposition}

\begin{proof}
(1) 在零组合中，考察出现的最高次数：其
系数必为零（没有别的项能抵消该次数），并
逐次向下递推。

(2) 若 $x \in \operatorname{Vect}(x_1, \dots, x_p)$，则关系式 $x
- \sum\lambda_i x_i = 0$ 是非平凡的。反之，
$(x_1, \dots, x_p, x)$ 的一个非平凡零组合必须
含有 $x$ 且其系数非零（否则与较小族的\omterm{def:b1:vspaces:free}{线性无关}性矛盾），
解出 $x$ 即知它位于\omterm{def:b1:vspaces:span}{张成的子空间}之中。

(3) 子族的零组合即为整个族的一个零组合，
其缺失项的系数取为 $0$；由
大族的\omterm{def:b1:vspaces:free}{线性无关}性使所有系数都为零。超族：
$E$ 中每个向量都已是生成部分的组合；
对多出的向量取系数为 $0$ 即可。
\end{proof}

\begin{example}[阶梯原理]
\label{ex:b1:vspaces:staircase}
设 $P_0, P_1, \dots, P_n \in K_n[X]$，且 $\deg P_k = k$（对每个
$k$，即次数构成“阶梯”）。则 $(P_0, \dots, P_n)$ 是
$K_n[X]$ 的一组基。\omterm{def:b1:vspaces:free}{线性无关}性
由\cref{prop:b1:vspaces:freecriteria}~(1)可知。关于生成性，
对次数用有限下降法论证：设 $Q \in
K_n[X]$，$Q \neq 0$，次数为 $d$，首项系数为 $a$，
并设 $b \neq 0$ 为 $P_d$ 的首项系数。则 $Q -
\frac ab P_d$ 的次数 $< d$（最高次项相消）；以此差替换
$Q$ 并反复进行，至多 $n + 1$ 步后得到零\omterm{def:b1:poly:def}{多项式}，
逆推这些减法，
便把 $Q$ 表示为 $P_k$ 的组合。此前已见过两个阶梯：
平移幂族 $\bigl((X-a)^k\bigr)_{0 \leq k
\leq n}$（见\cref{exo:b1:vspaces:4}），以及牛顿乘积
$\bigl((X - x_0)(X - x_1)\cdots(X - x_{k-1})\bigr)_{0 \leq k \leq
n}$，并在周末问题中加以运用。
\end{example}

\begin{example}[函数空间中的线性无关性]\label{ex:b1:vspaces:functionfree}
在 $\mathcal{F}(\R,\R)$ 中，族 $(\eu^{a_1 x}, \dots, \eu^{a_p
x})$（其中 $a_1 < \dots < a_p$）\omterm{def:b1:vspaces:free}{线性无关}：将零组合
除以 $\eu^{a_p x}$，并令 $x \to +\infty$；最后一个系数必为零，
然后逐次向下递推（见\cref{exo:b1:vspaces:8}，其中详述此论证及
其变体）。函数的\omterm{def:b1:vspaces:free}{线性无关}性靠\emph{求值}来证明：
在恰当的点处求值、在无穷远处考察，或先求导再求值。
\end{example}

\begin{example}[隐藏的关系使张成的子空间变小]
\label{ex:b1:vspaces:hiddenrelation}
在 $\mathcal{F}(\R, \R)$ 中，
$\operatorname{Vect}(1,\ \cos^2,\ \sin^2)$ 是什么？恒等式
$\cos^2 + \sin^2 = 1$ 是一个非平凡的零组合
\[
1\cdot\mathbf{1} + (-1)\cos^2 + (-1)\sin^2 = 0 :
\]
该族线性相关，而\omterm{def:b1:vspaces:span}{张成的子空间}仅由
$(1, \cos^2)$ 生成即可（$\sin^2 = 1 - \cos^2$）。该较小的族
\omterm{def:b1:vspaces:free}{线性无关}：若 $a + b\cos^2 x = 0$ 对一切 $x$ 成立，则取 $x = 0$ 与
$x = \frac\pi2$ 得 $a + b = 0$ 且 $a = 0$。于是\omterm{def:b1:vspaces:span}{张成的子空间}
是函数空间中的一个\emph{平面}——它还包含
$\cos 2x = 2\cos^2 x - 1$：形式上呈线性的三角函数族
在恒等式作用下常常退化，
这正是\omterm{def:b1:vspaces:free}{线性无关}性必须\emph{证明}，绝不能仅凭
族所含向量的个数而予以假定。
\end{example}

\begin{remark}[常见陷阱]
\label{rem:b1:vspaces:pitfalls}
四个经典陷阱。
\emph{逐对考察并不足够}：在 $\R^2$ 中，向量 $(1,0)$，
$(0,1)$，$(1,1)$ 两两不成比例，但却线性相关——
\omterm{def:b1:vspaces:free}{线性无关}是\emph{整个}族的性质，须由一个整体的零组合来检验，
绝非两两检验。
\emph{零向量会破坏一切}：任何包含
$0$ 的族都线性相关（$1\cdot 0 = 0$ 是一个非平凡的关系），
无论其余向量多么无辜。
\emph{并集不是和}：$F \cup G$ 几乎从不是\omterm{def:b1:vspaces:subspace}{子空间}
（见\cref{def:b1:vspaces:subspace}）；包含二者的最小\omterm{def:b1:vspaces:subspace}{子空间}
是 $F + G$，通常远大于并集——
在 $\R^2$ 中，两条不同直线的并是一个十字形，
而其和恰为整个平面。
\emph{直和要求交集平凡，而非不相交}：
两个\omterm{def:b1:vspaces:subspace}{子空间}永远不会不相交（二者都含 $0$）；正确的
条件是 $F \cap G = \{0\}$，且它必须加以\emph{证明}，
而不能凭图形直接读出——（参见
\cref{ex:b1:vspaces:manysupplements}，其中许多不同的 $G$
均可行）。
\emph{\omterm{def:b1:vspaces:free}{线性无关}性依赖于标量}：向量组 $(1, \iu)$
在 $\C$ 视为 $\R$-向量空间时\omterm{def:b1:vspaces:free}{线性无关}，而在 $\C$
视为 $\C$-向量空间时则线性相关（$\iu\cdot 1 + (-1)\cdot\iu = 0$）。
判定某族\omterm{def:b1:vspaces:free}{线性无关}之前，务必明确起作用的是哪一个域，
——\cref{ch:b1:findim}中的周末问题正是把
这种敏感性转化为无理性证明。
\end{remark}

\begin{remark}[这套语言的后续应用]
\label{rem:b1:vspaces:whereused}
本章之后的一切内容都在此建立的语言之上展开。
\cref{ch:b1:findim}计数基向量，把“\omterm{def:b1:vspaces:free}{线性无关}”与
“生成”转化为关于单个整数的不等式，该整数即
维数。\cref{ch:b1:linmaps}研究与这两种运算相容的\omterm{def:b1:logic:map}{映射}；
直和在那里化为投影。
\cref{ch:b1:matrices}用向量在某个
基下的\omterm{prop:b1:vspaces:coordinates}{坐标}对其编码——\cref{prop:b1:vspaces:coordinates}为这种编码提供依据，
而\cref{ch:b1:euclid}引入长度与夹角，
其建立在线性结构之上。在本科二年级分册中，同一组公理
公理（一字不改）适用于任意域，并可用于无穷
维空间；本章中任何地方都未用到有限性。
\end{remark}

\begin{remark}[贯穿第 3 册的三条线索]
\label{rem:b1:vspaces:threads}
留意本章三个具体思想如何逐步成长。
\emph{阶梯原理}
（\cref{ex:b1:vspaces:staircase}）在本章的周末问题中化为牛顿基，
在该周末问题中化为二项式基 $(B_k)$，
又在\cref{ch:b1:det}的周末问题中化为\omterm{def:b1:poly:def}{多项式}交错技巧：
一条引理，
三份无需行列式的收获。
\emph{求值作无关性检验}
（\cref{ex:b1:vspaces:functionfree}）化为\cref{ch:b1:linmaps}中的插值同构，
再化为\cref{ch:b1:det}中的范德蒙德判别法，
最后化为\cref{ch:b1:euclid}中的格拉姆检验：
同一种直觉，历经三次打磨。
\emph{直和}（\cref{def:b1:vspaces:sum}）化为
\cref{ch:b1:linmaps}中的投影，在\cref{ch:b1:euclid}中化为正交分解 $E =
F \oplus F^\perp$，以及
最小二乘法中“解释部分加残差”的分解，
在\cref{ch:b1:multivar}的周末问题中出现。本书极少内容，
归根结底，不过是这三种思想
换了新装而已。
\end{remark}

\section{练习}

\begin{exercise}[$\star$]\label{exo:b1:vspaces:1}
下列哪些是\omterm{def:b1:vspaces:subspace}{子空间}？
\begin{enumerate}
  \item $\{(x, y, z) \in \R^3 : x + 2y - z = 0\}$;
  \item $\{(x, y, z) \in \R^3 : x + 2y - z = 1\}$;
  \item $\{(x, y) \in \R^2 : xy \geq 0\}$;
  \item $\{P \in \R[X] : P(1) = 0\}$;
  \item $\{f \in \mathcal{F}(\R,\R) : f \text{ 有界}\}$。
\end{enumerate}
\end{exercise}

\begin{exercise}[$\star$]\label{exo:b1:vspaces:2}
在 $\R^3$ 中，$(1, 2, 1)$ 是否属于
$\operatorname{Vect}\bigl((1,0,1),\, (1,1,0)\bigr)$？还有 $(2, 1,
1)$ 呢？试用一个
方程刻画 $\operatorname{Vect}\bigl((1,0,1),(1,1,0)\bigr)$。
\end{exercise}

\begin{exercise}[$\star$]\label{exo:b1:vspaces:3}
判断 $\R^3$ 中下列向量组是否\omterm{def:b1:vspaces:free}{线性无关}：$\;\bigl((1,1,0), (1,0,1),
(0,1,1)\bigr)$; $\;\bigl((1,2,3), (2,4,6)\bigr)$;
$\;\bigl((1,0,0), (1,1,0), (1,1,1), (0,1,1)\bigr)$.
\end{exercise}

\begin{exercise}[$\star$]\label{exo:b1:vspaces:4}
证明 $(1, X - 1, (X-1)^2, (X-1)^3)$ 是 $\R_3[X]$ 的一组基，
并求 $X^3$ 在该基下的\omterm{prop:b1:vspaces:coordinates}{坐标}。\emph{（想想 $1$ 处的泰勒展开！）}
\end{exercise}

\begin{exercise}[$\star\star$]\label{exo:b1:vspaces:5}
在 $\R^4$ 中，设 $F = \{(x,y,z,t) : x = y = z\}$ 与 $G = \{(x,y,z,t)
: x = t = 0\}$。证明 $F \oplus G = \R^4$，并相应地
分解 $(1,2,3,4)$。
\end{exercise}

\begin{exercise}[$\star\star$]\label{exo:b1:vspaces:6}
在数列构成的空间中，设 $F$ 为全体收敛数列组成的
\omterm{def:b1:logic:sets}{集合}，$G = \operatorname{Vect}(u)$，其中 $u_n = (-1)^n$。
证明 $F \cap G = \{0\}$。$F + G$ 是否为
全体数列构成的空间？
\end{exercise}

\begin{exercise}[$\star\star$]\label{exo:b1:vspaces:7}
设 $F, G, H$ 为 $E$ 的\omterm{def:b1:vspaces:subspace}{子空间}。证明下式：
\[
F \cap (G + (F \cap H)) = (F \cap G) + (F \cap H),
\]
并在 $\R^2$ 中举例说明无限制的分配律
$F \cap (G + H) = (F\cap G) + (F \cap H)$ 不成立。
\end{exercise}

\begin{exercise}[$\star\star\star$]\label{exo:b1:vspaces:8}
证明 $\mathcal{F}(\R, \R)$ 中下列向量组\omterm{def:b1:vspaces:free}{线性无关}：
\begin{enumerate}
  \item $(\eu^{a_1 x}, \dots, \eu^{a_p x})$，其中 $a_1 < \dots < a_p$；
  \item $(\cos x, \sin x, \cos 2x, \sin 2x)$;
  \item $(x \mapsto \abs{x - a_1}, \dots, x \mapsto \abs{x - a_p})$
        其中 $a_i$ 互不相同 \emph{（每个函数
        恰在一点处不\omterm{def:b1:derivative:def}{可导}）}。
\end{enumerate}
\end{exercise}

\begin{exercise}[$\star\star\star$]\label{exo:b1:vspaces:9}
设 $E$ 为 $K$-向量空间，$F, G, H$ 为其\omterm{def:b1:vspaces:subspace}{子空间}，且满足 $F + G =
F + H$，$F \cap G = F \cap H$ 且 $G \subseteq H$。证明 $G = H$。
若去掉假设 $G \subseteq H$，给出一个反例。
\end{exercise}

\begin{exercise}[$\star\star$]\label{exo:b1:vspaces:10}
在 $\R[X]$ 中，设 $\mathcal P$ 为\emph{偶}\omterm{def:b1:poly:def}{多项式}构成的\omterm{def:b1:logic:sets}{集合}
（$P(-X) = P(X)$），且 $\mathcal I$ 为\emph{奇}\omterm{def:b1:poly:def}{多项式}构成的\omterm{def:b1:logic:sets}{集合}
（$P(-X) = -P(X)$）。证明 $\R[X] = \mathcal P \oplus \mathcal
I$，并证明 $\mathcal P =
\operatorname{Vect}(1, X^2, X^4, \dots)$，即偶
\omterm{def:b1:poly:def}{多项式}恰为关于 $X^2$ 的\omterm{def:b1:poly:def}{多项式}。
\end{exercise}

\begin{exercise}[$\star\star$]\label{exo:b1:vspaces:11}
设 $(x_1, x_2, x_3)$ 是实\omterm{def:b1:vspaces:def}{向量空间} $E$ 中的一个\omterm{def:b1:vspaces:free}{线性无关}向量组。
证明 $(x_1 + x_2,\ x_2 + x_3,\ x_3 + x_1)$ \omterm{def:b1:vspaces:free}{线性无关}。另问，
与之类似的四个向量组 $(x_1 + x_2,\ x_2 + x_3,\ x_3 +
x_4,\ x_4 + x_1)$ 在 $(x_1, x_2, x_3, x_4)$ \omterm{def:b1:vspaces:free}{线性无关}时是否也\omterm{def:b1:vspaces:free}{线性无关}？
\end{exercise}

\begin{exercise}[$\star\star\star$]\label{exo:b1:vspaces:12}
设 $E$ 为 $\R$（或 $\C$）上的\omterm{def:b1:vspaces:def}{向量空间}，$F_1, \dots,
F_k$ 为 $E$ 的\emph{真}\omterm{def:b1:vspaces:subspace}{子空间}（即每个 $F_i \neq E$）。
\begin{enumerate}
  \item 直接处理 $k = 2$ 的情形：若 $F_1 \not\subseteq
        F_2$ 且 $F_2 \not\subseteq F_1$，取 $x \in F_1
        \setminus F_2$，取 $y \in F_2 \setminus F_1$，并考察 $x
        + y$.
  \item 一般地证明 $E \neq F_1 \cup \dots \cup F_k$：无限域上的
        \omterm{def:b1:vspaces:def}{向量空间}永远不是有限个
        真\omterm{def:b1:vspaces:subspace}{子空间}的并。\emph{（取 $k$ 最小，
        取 $x \in F_1$ 使其不属于其余各 $F_i$，取 $y \notin F_1$，
        并沿直线 $t \mapsto y + tx$ 考察。）}
\end{enumerate}
\end{exercise}

\section{问题：插值，同一空间的三组基}

\begin{problem}\label{pb:b1:vspaces:1}
取定 $n + 1$ 个\emph{互异}的点 $x_0, x_1, \dots, x_n$，它们都属于 $\R$。
本题借助本章的观点，重新审视\omterm{thm:b1:poly:lagrange}{拉格朗日插值}
（见\cref{thm:b1:poly:lagrange}）。
空间 $\R_n[X]$ 上有三组自然的基——拉格朗日基、牛顿基，
以及（等距节点时的）二项式基——
每组基都使某一问题变得简单。本题的终点是一个真正的算术定理：
波利亚关于
把 $\Z$ 映入 $\Z$ 的\omterm{def:b1:poly:def}{多项式}的刻画。\index{拉格朗日插值}
\index{差商}\index{整值多项式}

\medskip
\noindent\textbf{第一部分——拉格朗日基。}对 $0 \leq i
\leq n$，令
\[
L_i \;=\; \prod_{j \neq i} \frac{X - x_j}{x_i - x_j} \;\in\;
\R_n[X].
\]
\begin{enumerate}
  \item 验证 $\deg L_i = n$，且 $L_i(x_j) = 1$（当 $j =
        i$ 时），$0$（当 $j \neq i$ 时）。
  \item 证明族 $(L_0, \dots, L_n)$ \omterm{def:b1:vspaces:free}{线性无关}。
  \item 证明：对每个 $P \in \R_n[X]$，
        \[
        P \;=\; \sum_{i=0}^{n} P(x_i)\, L_i ,
        \]
        由此得出 $(L_0, \dots, L_n)$ 是
        $\R_n[X]$ 的一组基。\emph{（考虑两端之差，
        并数其根的个数，见\cref{cor:b1:poly:nroots}。）}
  \item 推出插值定理：对任意 $y_0,
        \dots, y_n \in \R$，存在\emph{唯一}的 $P \in
        \R_n[X]$，使得 $P(x_i) = y_i$ 对所有 $i$ 都成立。在拉格朗日基下，
        \omterm{def:b1:poly:def}{多项式} $P$ 的\omterm{prop:b1:vspaces:coordinates}{坐标}是什么？
  \item 证明下列恒等式
        \[
        \sum_{i=0}^{n} L_i = 1
        \qquad\text{且对 } 0 \leq k \leq n,\qquad
        \sum_{i=0}^{n} x_i^{k}\, L_i = X^{k} .
        \]
\end{enumerate}

\medskip
\noindent\textbf{第二部分——牛顿基与
\omterm{pb:b1:vspaces:1}{差商}。}设 $N_0 = 1$，$N_k = (X - x_0)(X - x_1) \cdots
(X - x_{k-1})$，其中 $1 \leq k \leq n$。对定义在节点上的函数 $f$，
按如下方式定义\emph{\omterm{pb:b1:vspaces:1}{差商}}：
$f[x_i] = f(x_i)$，且
\[
f[x_i, \dots, x_{i+k}] \;=\;
\frac{f[x_{i+1}, \dots, x_{i+k}] - f[x_i, \dots, x_{i+k-1}]}
{x_{i+k} - x_i} .
\]
\begin{enumerate}[resume]
  \item 证明 $(N_0, N_1, \dots, N_n)$ 是
        $\R_n[X]$ 的一组基。
  \item 将 $f[x_0, x_1]$ 与 $f[x_0, x_1, x_2]$ 用
        $f$ 的取值表示，然后计算
        $f(x) = x^2$ 在任意三个节点上的全部\omterm{pb:b1:vspaces:1}{差商}。
  \item （艾特肯引理）设 $R$ 为 $f$ 在 $x_0, \dots,
        x_{n-1}$ 上的插值，$Q$ 为 $f$ 在 $x_1, \dots, x_n$ 上的插值，
        二者次数均 $\leq n - 1$。证明下式
        \[
        S \;=\; \frac{(X - x_0)\,Q - (X - x_n)\,R}{x_n - x_0}
        \]
        是 $f$ 在 $x_0, x_1, \dots, x_n$ 上的插值。
  \item 对节点个数用归纳法推出
        $X^{k}$ 的系数（在 $f$ 于 $x_0,
        \dots, x_k$ 上的插值\omterm{def:b1:poly:def}{多项式}中）恰为 $f[x_0, \dots, x_k]$。
  \item 证明\emph{牛顿插值公式}：
        $f$ 在 $x_0, \dots, x_n$ 上的插值\omterm{def:b1:poly:def}{多项式}为
        \[
        P \;=\; \sum_{k=0}^{n} f[x_0, \dots, x_k]\, N_k ,
        \]
        并导出闭式
        \[
        f[x_0, \dots, x_k] \;=\; \sum_{i=0}^{k}
        \frac{f(x_i)}{\prod_{j \neq i,\, j \leq k} (x_i - x_j)} ,
        \]
        这表明 $f[x_0, \dots, x_k]$ 与
        节点的排列次序无关。
\end{enumerate}

\medskip
\noindent\textbf{第三部分——等距节点：
差分算子。}此后节点取为 $0, 1, 2, \dots$，且对
\omterm{def:b1:poly:def}{多项式} $P$，令
\[
\Delta P(X) = P(X + 1) - P(X),
\qquad
B_k = \frac{X(X-1)\cdots(X-k+1)}{k!} \quad (B_0 = 1).
\]
\begin{enumerate}[resume]
  \item 证明：若 $\deg P = m \geq 1$ 且首项系数为
        $a$，则 $\deg \Delta P = m - 1$，首项系数为
        $m\,a$，且 $\Delta$ 把常数映为零。
  \item 证明 $(B_0, B_1, \dots, B_n)$ 是
        $\R_n[X]$ 的一组基，且 $\Delta B_k = B_{k-1}$（$k \geq 1$）。
  \item （牛顿前差公式）证明：对每个 $P
        \in \R_n[X]$ 满足下式
        \[
        P \;=\; \sum_{k=0}^{n} \bigl(\Delta^{k} P\bigr)(0)\, B_k .
        \]
  \item 证明：对每个 $k \geq 0$，
        \[
        \bigl(\Delta^{k} P\bigr)(0) \;=\;
        \sum_{j=0}^{k} (-1)^{k-j} \binom{k}{j} P(j) .
        \]
  \item 证明：若 $\deg P = n$ 且首项系数为 $a_n$，
        则 $\Delta^{n} P$ 为常数 $n!\,a_n$，且
        $\Delta^{n+1} P = 0$。
\end{enumerate}

\medskip
\noindent\textbf{第四部分——\omterm{pb:b1:vspaces:1}{整值多项式}。}称 $P \in \R[X]$ 为
\emph{\omterm{pb:b1:vspaces:1}{整值多项式}}，当 $P(m) \in
\Z$ 对所有 $m \in \Z$ 成立时。
\begin{enumerate}[resume]
  \item 证明每个 $B_k$ 都是\omterm{pb:b1:vspaces:1}{整值多项式}。\emph{（分别讨论 $m
        \geq k$、$0 \leq m < k$ 与 $m < 0$ 三种情形；当 $m =
        -q < 0$ 时，证明 $B_k(-q) = (-1)^k \binom{q + k - 1}{k}$。）}
  \item 证明\emph{波利亚刻画}：$P \in \R_n[X]$
        是\omterm{pb:b1:vspaces:1}{整值多项式}，当且仅当其关于基
        $(B_0, \dots, B_n)$ 的\omterm{prop:b1:vspaces:coordinates}{坐标}均为整数。
  \item 由此推出：若 $P \in \R_n[X]$ 在 $n + 1$
        个\emph{\omterm{def:b1:continuity:continuous}{连续}}整数 $a, a+1, \dots, a+n$ 处取整数值，则 $P$
        为\omterm{pb:b1:vspaces:1}{整值多项式}。\emph{（平移：对 $Q(X)
        = P(X + a)$ 应用上述讨论。）}
  \item 由第 16 题推出，$k$ 个
        \omterm{def:b1:continuity:continuous}{连续}整数之积总能被 $k!$ \omterm{def:b1:arith:divides}{整除}。
  \item 设 $P = \dfrac{X(X+1)(2X+1)}{6}$。求其在 $0, 1, 2, 3$ 处的
        牛顿表，将 $P$ 表示为基 $(B_k)$ 的组合，
        并由此断定：$P$ 是\omterm{pb:b1:vspaces:1}{整值多项式}，尽管其
        单项式系数中没有一个是整数。验证 $\Delta P =
        (X+1)^2$，并由此推出 $P(m) = 1^2 + 2^2 + \dots + m^2$，
        其中 $m \in \N$。
\end{enumerate}

\medskip
\noindent\textbf{第五部分——红利。}
\begin{enumerate}[resume]
  \item 设 $P \in \R_n[X]$ 以 $2^i$ 为插值数据，其中 $i =
        0, 1, \dots, n$。证明 $P = B_0 + B_1 + \dots + B_n$，且
        $P(n + 1) = 2^{n+1} - 1$：“倍增规律”
        总在紧接着的下一个点处失效。
  \item （离散原函数）证明：对所有整数 $m
        \geq 1$ 且 $k \geq 0$，
        \[
        \sum_{j=0}^{m-1} B_k(j) \;=\; B_{k+1}(m),
        \]
        即曲棍球棒恒等式 $\sum_{j=k}^{m-1}
        \binom{j}{k} = \binom{m}{k+1}$。
  \item 将 $X^2$ 与 $X^3$ 按基 $(B_k)$ 展开，并由此给出
        $\sum_{j=0}^{m-1} j^2$ 与
        $\sum_{j=0}^{m-1} j^3$ 的闭式，并由此导出尼科马库斯恒等式
        $1^3 + \dots + m^3 = (1 + \dots + m)^2$。
  \item 取 $n = 2$，节点为 $0, 1, 2$。写出
        $X^2$ 在本题三组基下的\omterm{prop:b1:vspaces:coordinates}{坐标}，依次为单项式基、
        拉格朗日基与牛顿基下的\omterm{prop:b1:vspaces:coordinates}{坐标}。检验
        三个答案与第 4 题和第 9 题是否相符。
  \item 综合题。用四句话回答：哪一个\omterm{def:b1:vspaces:def}{向量空间}概念使第 4 题
        自动成立；为何牛顿基\emph{递归地}计算\omterm{prop:b1:vspaces:coordinates}{坐标}，
        而拉格朗日基则\emph{即时地}读出\omterm{prop:b1:vspaces:coordinates}{坐标}；两组基共同满足的
        \omterm{def:b1:vspaces:free}{线性无关}判别准则是什么；
        并说明波利亚定理在何种确切意义下断言，\omterm{def:b1:poly:def}{多项式}的整值性
        是其在
        \emph{正确的基}下的\omterm{prop:b1:vspaces:coordinates}{坐标}的性质。
\end{enumerate}
\end{problem}
